% =============================================================================
% Technical note: validation of the sim_lighttrap Geant4 package
% JINST template (jinstpub.sty, https://jinst.sissa.it/jinst/help/JINST_TeXclass.jsp)
%
% Build:  pdflatex lighttrap_validation_note
%         pdflatex lighttrap_validation_note      % twice, for cross-references
%
% INTERNAL RECORD ONLY -- not intended for publication.
% =============================================================================
\documentclass[a4paper,11pt]{article}
\pdfoutput=1

\usepackage{jinstpub}

\usepackage{booktabs}
\usepackage{amsmath}
\usepackage{siunitx}
\usepackage{xcolor}
\usepackage{listings}

\lstset{
  basicstyle=\ttfamily\footnotesize,
  breaklines=true,
  frame=single,
  framesep=4pt,
  rulecolor=\color{gray!50},
  backgroundcolor=\color{gray!5},
  columns=fullflexible,
  keepspaces=true,
}

\title{\boldmath Validation of the \texttt{sim\_lighttrap} Geant4 optical
simulation package}

\author[a,1]{R.~Razakamiandra}
\author[a]{and W.~Shi}

\affiliation[a]{Department of Physics and Astronomy, Stony Brook University,\\
Stony Brook, NY 11794, USA}

\emailAdd{radofanantenan.razakamiandra@stonybrook.edu}

\note[1]{Corresponding author.}

\abstract{We report a code-level and run-level validation of
\texttt{sim\_lighttrap}, a Geant4 application that models a two-stage
wavelength-shifting (WLS) light-trap module immersed in liquid argon. The
package was configured and compiled from a clean tree against Geant4~11.3.2
(installed via Spack) and ROOT~6.40.02 on macOS/arm64, and exercised with its
unit-test suite and with a set of controlled photon-transport runs. The build
is clean and the source-level test suite passes in full. The photon-transport
results, however, do not reproduce the behaviour the package documentation
predicts for the selectable backplane reflector: the module collects
$3.3\times$ \emph{fewer} photoelectrons with a reflective backplane than with
no backplane at all. We trace this to three specific, independently verified
properties of the current optical model --- an open (unmirrored, non-dichroic)
front face on the blue-WLS slab, an optically coupled \texttt{dielectric\_metal}
skin surface that suppresses total internal reflection at the back face, and a
surface-finish/type combination that renders the nominal ``diffuse PTFE''
option bit-for-bit identical to the specular Vikuiti option. We document the
evidence for each, quantify the resulting efficiency chain, and list the
follow-up items required before the package is used for detector-performance
predictions. This note is an internal record; the findings are stated for human
review and have not been independently confirmed.}

\keywords{Simulation methods and programs; Photon detectors for UV, visible and
IR photons; Noble liquid detectors; Detector modelling and simulations}

\begin{document}
\maketitle
\flushbottom

% =============================================================================
\section{Scope and status}
\label{sec:scope}

This note records an independent read-through and execution-level validation of
the \texttt{sim\_lighttrap} package at commit \texttt{389e967} (branch
\texttt{develop}, 2026-07-21). It is an internal engineering record, not a
physics result and not a publication. Two classes of statement appear below and
should be read differently:

\begin{itemize}
\item \textbf{Verified.} Statements backed by a command that was executed and
whose output is reproduced or summarised here. These include the build, the
unit-test result, and every number in
sections~\ref{sec:results}--\ref{sec:findings}.
\item \textbf{Interpretation.} Statements about \emph{why} the simulation
behaves as observed. These are supported by targeted experiments described in
the text, but they are the author's reading of the evidence and require
confirmation by a human reviewer familiar with the intended detector design.
\end{itemize}

No source file of the simulation was modified during this validation. All runs
were performed in a scratch directory outside the repository working tree.

% =============================================================================
\section{Software environment}
\label{sec:env}

The host and toolchain used for the validation are listed in
table~\ref{tab:env}. Geant4 is provided by a Spack installation; the relevant
variants are recorded because several of them (notably \texttt{cxxstd=17} and
the absence of \texttt{+tbb}) affect how the application is built and whether
multithreaded event processing is available.

\begin{table}[htbp]
\centering
\caption{\label{tab:env}Validation environment.}
\begin{tabular}{ll}
\toprule
Component & Version / detail \\
\midrule
Host            & macOS (Darwin 25.6.0), Apple silicon (arm64) \\
Geant4          & 11.3.2, installed with Spack \\
Geant4 spec     & \texttt{geant4@11.3.2 +data +opengl +qt +threads} \\
                & \texttt{\textasciitilde{}hdf5 \textasciitilde{}tbb \textasciitilde{}vecgeom \textasciitilde{}vtk \textasciitilde{}x11}, \texttt{cxxstd=17}, \texttt{build\_type=Release} \\
Geant4 prefix   & \texttt{\$SPACK\_ROOT/opt/spack/darwin-m1/geant4-11.3.2-3saow3g...} \\
ROOT            & 6.40.02 (Homebrew), used for output inspection and analysis \\
Compiler        & Apple Clang, \texttt{/usr/bin/c++} \\
CMake           & 4.x (Homebrew) \\
yaml-cpp        & resolved by \texttt{find\_package}; \texttt{FetchContent} fallback available \\
Python          & 3.13, \texttt{pytest} 9.1.1 \\
\bottomrule
\end{tabular}
\end{table}

\paragraph{Geant4 discovery.} \texttt{CMakeLists.txt} deliberately avoids a
hard-coded Geant4 path and accepts, in order of precedence, \texttt{-DGeant4\_DIR},
\texttt{-DCMAKE\_PREFIX\_PATH}, and the \texttt{GEANT4\_CMAKE\_DIR} environment
variable. This is the correct pattern for a Spack-provided Geant4 and worked
without modification.

\paragraph{Toolchain caveat.} \label{par:conda} On the validation host a
Miniconda \texttt{clang++} precedes the system compiler on \texttt{PATH}, and
CMake selects it by default. Its linker driver fails the compiler-identification
step with

\begin{lstlisting}
ld: -lto_library library filename must be 'libLTO.dylib'
\end{lstlisting}

which aborts configuration before Geant4 is even searched for. The failure is an
environment artefact, not a defect of the package. Configuring with the system
compiler explicitly avoids it, and this is the invocation used throughout:

\begin{lstlisting}
cmake -DCMAKE_C_COMPILER=/usr/bin/cc \
      -DCMAKE_CXX_COMPILER=/usr/bin/c++ \
      -DGeant4_DIR=$G4_PREFIX/lib/cmake/Geant4 \
      /path/to/sim_lighttrap
make -j8
\end{lstlisting}

The existing in-tree \texttt{build/CMakeCache.txt} already records
\texttt{/usr/bin/c++}, so the working build was configured the same way. It is
worth pinning the compiler in the project documentation.

% =============================================================================
\section{Package overview}
\label{sec:package}

\subsection{Layout}

The application follows the conventional Geant4 example structure:
\texttt{src/} and \texttt{include/} hold the C++ sources
($\sim$1\,500 lines of implementation), \texttt{mac/} the Geant4 macros,
\texttt{input/} the tabulated data files (SiPM photon-detection efficiency, WLS
emission spectra), \texttt{analysis/} the ROOT/PyROOT post-processing,
\texttt{tests/} a \texttt{pytest} suite, and \texttt{docs/} Org-mode design
notes. A single YAML file, \texttt{sim\_config.yaml}, exposes the full set of
tunable parameters.

\subsection{Configuration system}

Configuration is layered, which is a genuine strength of the package. Geometry,
source and run parameters may be supplied by YAML, by Geant4 macro commands, or
by both:

\begin{lstlisting}
./sim_lighttrap sim_config.yaml            # YAML only
./sim_lighttrap run.mac                    # macro only
./sim_lighttrap sim_config.yaml run.mac    # YAML first, macro overrides
\end{lstlisting}

\texttt{src/sim\_lighttrap.cc} distinguishes the two file types by extension,
parses the YAML into a \texttt{SimConfig} struct (\texttt{src/config.cc}), and
then \emph{replays} it through the same \texttt{G4UImanager} commands a macro
would use (\texttt{/detector/...}, \texttt{/gps/...}), followed by
\texttt{/run/reinitializeGeometry}. Routing YAML through the messenger rather
than through a second code path means both entry points exercise identical
logic, which is good practice.

Per-material optical properties are also YAML-configurable on a shared 13-point
wavelength grid $\{530, 425, 400, 340, 305, 160, 145, 135, 128, 125, 120, 115,
106\}$\,nm. \texttt{validateOpticalArray()} (\texttt{include/config.hh}) enforces
the 13-entry length at parse time and throws otherwise; any absent key falls
back to a hard-coded default in \texttt{construction.cc}. The bounds guarantee
this provides is what makes the unchecked fixed-size copies in
\texttt{resolveArray()} and \texttt{addWLSEmission()} safe.

WLS emission spectra additionally accept an arbitrary-length two-column file
(\texttt{wlscomponent\_file}), which is converted to strictly ascending photon
energy with duplicates dropped before being handed to Geant4 --- a correct
handling of a requirement that is easy to get wrong. Precedence is
file~$>$~inline 13-point~$>$~built-in smooth asymmetric Gaussian. In the
as-shipped configuration both file entries are commented out, so the runs
reported here use the built-in Gaussian defaults (pTP peak 340\,nm, blue-WLS
peak 430\,nm).

\subsection{Geometry}
\label{sec:geometry}

The module is built in \texttt{MyLightTrapConstruction::Construct()}
(\texttt{src/construction.cc}). All $z$-coordinates are derived from a single
reference --- the back face of the blue-WLS slab at $z=0$ --- so that changing
any thickness shifts every downstream component consistently. This is a clean
design and was verified by inspection. Photons travel in $+z$. The default stack
is given in table~\ref{tab:geom}.

\begin{table}[htbp]
\centering
\caption{\label{tab:geom}Default module stack. The $z$ range is given in the
convention $z=0$ at the blue-WLS back face.}
\begin{tabular}{llll}
\toprule
Component & Material & Thickness & $z$ range [mm] \\
\midrule
pTP WLS film            & \texttt{pTP} (C$_{18}$H$_{14}$)     & \SI{2}{\micro\meter} & $-12.002$ to $-12$ \\
UV-transparent acrylic  & \texttt{uvTransAcrylic}             & \SI{3}{\milli\meter} & $-12$ to $-9$ \\
LAr gap                 & \texttt{G4\_lAr} (world)            & \SI{3}{\milli\meter} & $-9$ to $-6$ \\
Blue-WLS slab           & \texttt{bluewlsacrylic}             & \SI{6}{\milli\meter} & $-6$ to $0$ \\
Backplane foil          & \texttt{acrylicMcMaster} + skin     & \SI{65}{\micro\meter} & $0$ to $0.065$ \\
Backplane leak counter  & \texttt{G4\_lAr}                    & \SI{0.1}{\milli\meter} & $0.065$ to $0.165$ \\
\bottomrule
\end{tabular}
\end{table}

The module face is $15\times15$\,cm$^2$. Thirty SiPMs, each presenting a
$6\times6$\,mm$^2$ face, are placed flush against the $\pm y$ edges of the
blue-WLS slab (15 per edge, uniform pitch). Vikuiti foils cover the $\pm x$
edges. The world is a liquid-argon box 60\,\% larger than the module on each
side.

Three additional volumes exist purely as instrumentation. They are built from
the world material (LAr), so they are refractive-index-matched to their
surroundings and introduce no optical boundary; they only count photons. The
first-layer backplane counter is offset \SI{0.01}{\milli\meter} behind the
first-layer back face so that totally internally reflected photons are not
counted --- a careful and correct detail. Two of the three counters are
pass-through (they do not kill the track); the backplane leak counter is
terminal.

Overlap checking is enabled on every placement (\texttt{checkOverlaps=true}) and
reported \texttt{OK} for all volumes in every configuration run.

\subsection{Physics and primary generation}

The physics list (\texttt{src/physics.cc}) registers
\texttt{G4EmStandardPhysics} and \texttt{G4OpticalPhysics} only. For a study
whose primaries are optical photons this is adequate, though it means the
package cannot currently be driven by ionising radiation depositing energy in
the argon.

Primaries are generated with \texttt{G4GeneralParticleSource}
(\texttt{src/generator.cc}). The default run fires \SI{9.68}{\electronvolt}
(\SI{128}{\nano\meter}) optical photons from a $15\times15$\,cm$^2$ planar
square source at $z=\SI{-2}{\centi\meter}$, directed along $+z$.

\subsection{Output data model}

\texttt{MyRunAction} (\texttt{src/run.cc}) defines nine ntuples, written to
\texttt{output\_apexltsim<runID>.root}. They are summarised in
table~\ref{tab:ntuples}.

\begin{table}[htbp]
\centering
\caption{\label{tab:ntuples}Output ntuples.}
\begin{tabular}{cll}
\toprule
ID & Name & Content \\
\midrule
0 & \texttt{SiPMPhotons}          & every photon reaching a SiPM face (before PDE) \\
1 & \texttt{SiPMHits}             & subset passing the per-wavelength PDE roll \\
2 & \texttt{Energy}               & deposited energy per event \\
3 & \texttt{Primary}              & primary particle kinematics \\
4 & \texttt{WLSPhoton}            & WLS / scintillation photon births \\
5 & \texttt{OpticalValidation}    & primary reflection / incidence bookkeeping \\
6 & \texttt{SecondLayerBackplane} & photons reaching the backplane leak counter \\
7 & \texttt{FirstLayerBackplane}  & crossings of the first-layer back face \\
8 & \texttt{FirstLayerEdges}      & crossings at the first-layer $\pm x$ edges \\
\bottomrule
\end{tabular}
\end{table}

Ntuples 6--8 share a common fill in \texttt{MyLeakDetector::ProcessHits} and
carry a \texttt{forward} flag ($1$ if $p_z>0$) alongside the direction cosine,
so forward and back-reflected traffic can be separated in analysis. This flag is
used extensively in section~\ref{sec:findings}. Detected hits are produced by
rolling against the tabulated Broadcom PDE curve in
\texttt{input/pde\_broadcom.dat} (401 points, peak $0.63$ at \SI{420}{\nano\meter}).

% =============================================================================
\section{Validation procedure and results}
\label{sec:results}

\subsection{Build}

Configuration and compilation from a clean out-of-source build directory
succeeded with the invocation of section~\ref{sec:env}: \textbf{exit status 0,
zero compiler warnings, zero errors}, producing the \texttt{sim\_lighttrap}
executable. CMake emits one warning originating inside the Qt6 dependency chain
pulled in by \texttt{Geant4Config.cmake} (a request to link the directory
\texttt{/usr/X11R6/lib}, which CMake drops). It is upstream and harmless.

\subsection{Unit tests}

The \texttt{pytest} suite is a source-level suite: it inspects the C++ and YAML
sources for expected declarations, parameter names and documentation, and
requires neither a build nor a simulation run. It completed as

\begin{lstlisting}
319 passed, 6 skipped in 2.77s
\end{lstlisting}

This confirms that the source text is internally consistent --- the messenger
commands, YAML keys, ntuple definitions and README statements agree with one
another. It is worth being explicit, however, that no test in this suite
executes the simulation or asserts anything about photon transport. The suite
therefore cannot detect the behaviour reported in section~\ref{sec:findings}, and
its passing should not be read as physics validation.

\subsection{Transport runs}

Four configurations were run, each with $2\times10^{4}$ primary photons and
otherwise the shipped defaults. Three use the selectable backplane reflector
options; the fourth is a diagnostic described in
section~\ref{sec:findings}. Results are given in table~\ref{tab:eff}, computed
with the package's own \texttt{analysis/analyze\_efficiency.py}.

\begin{table}[htbp]
\centering
\caption{\label{tab:eff}Efficiency chain for the three backplane options,
$N_{\text{primary}}=2\times10^{4}$ VUV photons at \SI{128}{\nano\meter}.
Uncertainties are the binomial errors reported by the package's analysis script.}
\begin{tabular}{lccc}
\toprule
Quantity & \texttt{vikuiti} & \texttt{ptfe} & \texttt{none} \\
\midrule
pTP WLS photons $N_{\text{pTP}}$        & 19410 & 19415 & 19464 \\
$\varepsilon_{\text{WLS1}}=N_{\text{pTP}}/N_{\text{primary}}$
  & $97.05\pm0.12\,\%$ & $97.08\pm0.12\,\%$ & $97.32\pm0.11\,\%$ \\
Blue-WLS photons $N_{\text{blue}}$      & 10658 & 10598 & 10938 \\
$N_{\text{blue}}/N_{\text{pTP}}$
  & $54.91\pm0.36\,\%$ & $54.59\pm0.36\,\%$ & $56.20\pm0.36\,\%$ \\
\midrule
Photons at SiPM face $N_{\text{at SiPM}}$ & 1019 & 1014 & \textbf{3493} \\
Collection $N_{\text{at SiPM}}/N_{\text{blue}}$
  & $9.56\pm0.29\,\%$ & $9.57\pm0.29\,\%$ & $\mathbf{31.94\pm0.45\,\%}$ \\
Detected $N_{\text{hits}}$              & 616 & 601 & \textbf{2046} \\
Detection $N_{\text{hits}}/N_{\text{blue}}$
  & $5.78\pm0.23\,\%$ & $5.67\pm0.23\,\%$ & $\mathbf{18.71\pm0.37\,\%}$ \\
Average PDE $N_{\text{hits}}/N_{\text{at SiPM}}$
  & $60.5\pm1.5\,\%$ & $59.3\pm1.5\,\%$ & $58.6\pm0.8\,\%$ \\
\midrule
\textbf{p.e.\ per primary} $N_{\text{hits}}/N_{\text{primary}}$
  & $\mathbf{3.08\pm0.12\,\%}$ & $\mathbf{3.01\pm0.12\,\%}$ & $\mathbf{10.23\pm0.21\,\%}$ \\
\midrule
\texttt{SecondLayerBackplane} (ntuple 6) & 0 & 1 & 3028 \\
\texttt{FirstLayerEdges} (ntuple 8)      & 579 & 581 & 510 \\
\bottomrule
\end{tabular}
\end{table}

Two features of table~\ref{tab:eff} are behaving as designed and are worth
recording as positive results. First, the first-stage conversion efficiency
$\varepsilon_{\text{WLS1}}\approx97\,\%$ is stable across all configurations, as
it must be: the backplane is downstream of the pTP film and cannot affect it.
Second, the average realised PDE, $58$--$60\,\%$, is consistent with the
Broadcom curve peak of $63\,\%$ folded against a $\sim$\SI{430}{\nano\meter}
emission spectrum. The first-stage optical model and the PDE application are
therefore self-consistent.

The remaining entries are not as designed, and are the subject of the next
section.

% =============================================================================
\section{Findings}
\label{sec:findings}

\subsection{F1: the backplane reflector reduces light collection by \texorpdfstring{$3.3\times$}{3.3x}}
\label{sec:f1}

The package README states the expected ordering of the three backplane options
explicitly:

\begin{quote}
``\emph{Collection / detection efficiency (\texttt{SiPMHits}, ntuple 1) ---
typically \texttt{vikuiti} $\geq$ \texttt{ptfe} $>$ \texttt{none}.}''
\end{quote}

The simulation produces the opposite ordering, by a large margin. From
table~\ref{tab:eff}, removing the reflector raises the photoelectron yield per
primary from $3.08\,\%$ to $10.23\,\%$, a factor $3.3$. The separation is far
outside the quoted uncertainties. Adding a nominally $98\,\%$ reflective
backplane to this geometry makes the module substantially \emph{worse}.

\paragraph{Evidence for the mechanism.} Three experiments localise the cause.

\emph{(i) It is not absorption at the foil.} Setting the Vikuiti reflectivity to
$1.0$ at every grid point --- a lossless mirror --- gives $N_{\text{hits}}=636$,
statistically indistinguishable from the $616$ obtained at $R=0.98$ and still
far below the $2046$ obtained with no foil. Per-bounce absorption in the foil is
therefore not the loss channel.

\emph{(ii) The lost photons are blue, and they leave through the front of the
module.} The wavelength spectrum at the SiPM face shows that the deficit is
confined entirely to the \SI{400}{}--\SI{530}{\nano\meter} band emitted by the
blue-WLS slab; the \SI{305}{}--\SI{400}{\nano\meter} pTP component is unchanged
between configurations. Splitting the first-layer backplane counter (ntuple 7)
by the \texttt{forward} flag and by wavelength locates where that blue light
goes (table~\ref{tab:leak}).

\begin{table}[htbp]
\centering
\caption{\label{tab:leak}Crossings of the first-layer backplane counter
(ntuple 7), split by direction and wavelength. ``Backward, blue'' counts blue
photons travelling in $-z$, i.e.\ out of the front of the module. These are
crossings, not unique photons, but the forward blue rate is small enough that
multiple counting is a minor correction.}
\begin{tabular}{lccc}
\toprule
Selection & \texttt{vikuiti} & \texttt{ptfe} & \texttt{none} \\
\midrule
Forward,  UV   ($\lambda<\SI{400}{\nano\meter}$) & 10556 & 10550 & 10588 \\
Forward,  blue ($\lambda>\SI{400}{\nano\meter}$) & 1152  & 1122  & 574 \\
Backward, UV   ($\lambda<\SI{400}{\nano\meter}$) & 528   & 538   & 496 \\
\textbf{Backward, blue} ($\lambda>\SI{400}{\nano\meter}$)
                                                 & \textbf{8805} & \textbf{8414} & \textbf{2521} \\
\midrule
Backward blue / $N_{\text{blue}}$                & $83\,\%$ & $79\,\%$ & $23\,\%$ \\
\bottomrule
\end{tabular}
\end{table}

With the reflector in place, $\sim$\,$83\,\%$ of all blue photons produced in
the slab exit through the \emph{front} of the module and are lost to the world
volume. With no reflector this falls to $23\,\%$; the rest either escape through
the back --- where they are counted by the leak detector, 3028 entries --- or
reach the SiPMs.

\emph{(iii) Interpretation.} The blue-WLS slab is a planar waveguide whose
guided modes are defined by total internal reflection at the
slab($n=1.58$)/LAr($n\approx1.23$) interfaces, giving a critical angle of
$\approx51^{\circ}$. The back face carries a \texttt{G4LogicalSkinSurface} of
type \texttt{dielectric\_metal}, placed flush against the slab. That surface
\emph{replaces} the dielectric interface: a photon arriving at the back face no
longer undergoes total internal reflection but a metallic reflection, which is
specular and therefore preserves the polar angle with respect to $z$. A ray
outside the guiding cone stays outside it. The mirror consequently cannot
convert an escaping mode into a trapped mode; it can only send the photon across
the slab to the opposite face --- which, critically, is \emph{open}. There is no
reflector and no dichroic filter on the front ($-z$) face of the blue-WLS slab
anywhere in the geometry. The reflector thus converts a back-face loss, which
the leak counter at least records, into a front-face loss out of the detector,
and additionally gives each photon repeated opportunities to find the open face.

This reading is consistent with every measurement above, but it is an
interpretation and should be confirmed by a reviewer who knows the intended
optical design. In particular, if the physical device is meant to include a
dichroic filter on its entrance window --- transmitting the \SI{340}{}--\SI{430}{\nano\meter}
pTP light inward while reflecting the \SI{430}{\nano\meter} blue light back in,
which is the defining element of a light trap of this class --- then that
component is simply absent from the geometry, and its absence is sufficient to
explain the entire effect. No occurrence of a dichroic or filter volume exists
in \texttt{src/}, \texttt{include/}, \texttt{sim\_config.yaml} or the Org-mode
design notes.

\subsection{F2: the PTFE option is bit-for-bit identical to the Vikuiti option}
\label{sec:f2}

The package presents the backplane choice as a specular-versus-diffuse
comparison. \texttt{DefinePTFEOpticalSurface()} sets

\begin{lstlisting}
PTFE->SetType(dielectric_metal);
PTFE->SetFinish(groundfrontpainted);   // "Lambertian diffuse reflection"
PTFE->SetModel(unified);
\end{lstlisting}

In Geant4's \texttt{G4OpBoundaryProcess}, \texttt{groundfrontpainted} selects
Lambertian reflection only on the \texttt{dielectric\_dielectric} branch. On the
\texttt{dielectric\_metal} branch, only \texttt{ground} and
\texttt{groundbackpainted} sample a perturbed facet normal; every other finish,
including \texttt{groundfrontpainted}, reflects about the unperturbed surface
normal, i.e.\ specularly. The combination as written therefore yields a specular
mirror, not a diffuse one.

\paragraph{Evidence.} The simulation uses a fixed random seed and no
randomisation (section~\ref{sec:repro}), so two configurations that differ only
in a label produce bit-identical output. Raising the PTFE reflectivity from
$0.95$ to $0.98$ --- making its only remaining difference from Vikuiti disappear
--- reproduces the Vikuiti run exactly (table~\ref{tab:ptfe}):

\begin{table}[htbp]
\centering
\caption{\label{tab:ptfe}PTFE at $R=0.98$ versus Vikuiti at $R=0.98$, same seed.}
\begin{tabular}{lcc}
\toprule
Ntuple & \texttt{vikuiti}, $R=0.98$ & \texttt{ptfe}, $R=0.98$ \\
\midrule
\texttt{SiPMPhotons}          & 1019  & 1019 \\
\texttt{SiPMHits}             & 616   & 616 \\
\texttt{SecondLayerBackplane} & 0     & 0 \\
\texttt{FirstLayerBackplane}  & 21041 & 21041 \\
\texttt{FirstLayerEdges}      & 579   & 579 \\
\bottomrule
\end{tabular}
\end{table}

Every counter agrees exactly. The two options differ only in the reflectivity
number; the specular-versus-diffuse comparison the package advertises is not
being performed. Note that this also means F1 cannot be rescued by the diffuse
option: a genuinely Lambertian backplane \emph{would} randomise the polar angle
and could feed photons back into guided modes, so it is a physically meaningful
thing to test --- but it has not yet been tested.

\subsection{F3: the backplane leak counter reads zero whenever a foil is present}
\label{sec:f3}

The README and the in-code comment describe ntuple 6 as recording ``every photon
transmitted through the backplane ($\sim2\,\%$ of those hitting it)''. It records
none: table~\ref{tab:eff} gives 0 entries for \texttt{vikuiti} and 1 for
\texttt{ptfe}, against an expectation of several hundred if $2$--$5\,\%$ of the
$\sim10^{4}$ incident photons were transmitted.

This is correct Geant4 behaviour and incorrect documentation. A
\texttt{dielectric\_metal} surface has two outcomes only: reflection with
probability $R$, or absorption with probability $1-R$. It never transmits. The
$(1-R)$ fraction is destroyed at the surface and never reaches the counter
\SI{65}{\micro\meter} behind it. The counter functions correctly in the
\texttt{none} configuration (3028 entries), where there is no skin surface and
the photons genuinely cross a dielectric interface.

The practical consequence is that the back-face escape observable is unavailable
for exactly the configurations in which it would be used, and that any analysis
comparing \texttt{SecondLayerBackplane} across foil options --- which the README
suggests doing --- compares a real measurement against a structural zero.

\subsection{F4: primary photon polarisation is set one event late}
\label{sec:f4}

Every run emits the Geant4 warning \texttt{G4Exception : ZeroPolarization}.
In \texttt{MyPrimaryGenerator::GeneratePrimaries}
(\texttt{src/generator.cc}) the vertex is created first and the polarisation is
assigned to the source afterwards:

\begin{lstlisting}
fGPS->GeneratePrimaryVertex(anEvent);        // vertex created here
...
fGPS->GetCurrentSource()->SetParticlePolarization(randomPol);  // applied later
\end{lstlisting}

The polarisation set at the end of event $n$ is therefore the one carried by the
primary of event $n+1$; the first event is generated with zero polarisation, and
every subsequent event uses the previous event's vector. Swapping the two
statements resolves both the warning and the off-by-one.

Separately, the sampling range is a quarter turn rather than a full one:

\begin{lstlisting}
// G4double phi = G4UniformRand() * 2.0 * CLHEP::pi;   // full range, commented out
G4double phi = G4UniformRand() * 0.5 * CLHEP::pi;      // active
\end{lstlisting}

so the polarisation azimuth is confined to $[0,\pi/2)$ and the primaries are
partially polarised rather than unpolarised. For the unpolarised-source study
intended here the full $2\pi$ range is the appropriate choice. The impact on the
present results is expected to be modest, since the dominant signal is
wavelength-shifted light whose polarisation is reset at emission, but it affects
the Fresnel treatment of the primaries at the first interface and should be
corrected before any surface-reflectivity study.

\subsection{F5: minor observations}
\label{sec:f5}

\begin{description}
\item[No bulk absorption in the blue-WLS slab.]
\texttt{DefineBlueWLSMaterial()} sets \texttt{RINDEX} and \texttt{WLSABSLENGTH}
but no \texttt{ABSLENGTH}; \texttt{DefineAcrylicMaterial()} sets only
\texttt{RINDEX}. Geant4 treats a missing \texttt{ABSLENGTH} as infinite, so the
slab and the foil bodies are bulk-lossless. Given that the intended geometry
guides light over $\sim$\SI{15}{\centi\meter} with many bounces, this is an
optimistic assumption and should be replaced with a measured attenuation length.

\item[Fixed output filename in the current working directory.]
\texttt{MyRunAction::BeginOfRunAction} writes
\texttt{output\_apexltsim<runID>.root} to the process working directory, with
run~0 always producing the same name. A previous long production file in
\texttt{build/} was destroyed this way. All runs for this note were performed in
a scratch directory for that reason. A configurable output path would remove a
standing hazard.

\item[Global mutation of a NIST material.] \texttt{DefineWorldMaterial()}
retrieves \texttt{G4\_lAr} from the NIST manager and attaches a properties table
to it, mutating the shared instance rather than a private copy. Harmless in this
single-geometry application; a latent problem if a second argon volume with
different properties is ever added.

\item[Deferred ntuple row in \texttt{OpticalValidation}.]
\texttt{src/tracking.cc} fills columns 0--3 and 5--7 in
\texttt{PreUserTrackingAction} and column 4 plus \texttt{AddNtupleRow} in
\texttt{PostUserTrackingAction}. This is correct in sequential mode, which the
application uses (\texttt{G4RunManager}), but relies on per-thread ntuple state
persisting across the track and would need review before enabling
multithreading.

\item[\texttt{FirstLayerBackplane} exceeds the primary count.] 21041 entries for
$2\times10^{4}$ primaries is expected and not a defect: the counter is
pass-through, so a single photon may cross it more than once, and the
\texttt{forward} flag is provided precisely so analyses can separate the two
directions. It is noted here only because the raw count invites
misinterpretation.
\end{description}

% =============================================================================
\section{Reproducibility}
\label{sec:repro}

No seed is set anywhere in \texttt{src/}, \texttt{include/} or \texttt{mac/}, so
every run starts from the Geant4 default engine state. Running the same
configuration twice gave $N_{\text{hits}}=616$ both times --- identical to the
last event.

This has a useful consequence and a limiting one. It makes the comparisons in
this note exact: the three backplane configurations differ only in the geometry
under test, with no statistical fluctuation between them, which is what allows
the bit-identity argument of section~\ref{sec:f2}. But it also means statistics
cannot be accumulated by repeating a run, and that the uncertainties quoted in
table~\ref{tab:eff} describe a single sample rather than a spread over
independent samples. Adding a \texttt{/random/setSeeds} hook, exposed through
the YAML configuration, would let both modes be used deliberately.

The full set of commands used for this validation is recorded in
\texttt{docs/technical\_note/README.md} alongside this source.

% =============================================================================
\section{Summary and recommended follow-up}
\label{sec:summary}

The engineering state of \texttt{sim\_lighttrap} is good. It builds cleanly and
without warnings against a Spack-provided Geant4~11.3.2, its layered
YAML/macro configuration routes both entry points through the same messenger
code, its optical-property tables are validated at parse time, its geometry is
parameterised from a single reference coordinate with overlap checking enabled
throughout, and its 319-test source-level suite passes in full. The first WLS
stage and the PDE application are self-consistent and behave as intended.

The optical model of the second stage is not yet giving trustworthy
detector-performance numbers. In the shipped configuration the module collects
$3.3\times$ fewer photoelectrons with its reflective backplane than without it,
the reverse of the documented expectation. The evidence assembled here points to
an open front face on the blue-WLS slab --- no reflector and, in particular, no
dichroic filter --- combined with an optically coupled \texttt{dielectric\_metal}
back surface that suppresses total internal reflection where the waveguide
relies on it. Two of the three diagnostics intended to detect exactly this class
of problem are themselves compromised: the diffuse-reflector option is not
diffuse, and the back-face leak counter is structurally zero whenever a foil is
present.

Recommended follow-up, in priority order:

\begin{enumerate}
\item \textbf{Confirm the intended optical design of the front face.} If a
dichroic filter window is part of the device, add it to the geometry; F1 is
expected to resolve with it in place. If the front face is genuinely open, then
the backplane reflector is counterproductive by construction and the README's
expected ordering should be corrected. This decision gates everything else.
\item \textbf{Fix the PTFE surface} (F2) --- either \texttt{SetType(dielectric\_dielectric)}
with \texttt{groundfrontpainted}, or \texttt{dielectric\_metal} with
\texttt{ground} and a non-zero \texttt{SigmaAlpha} --- so that a genuine
specular-versus-diffuse comparison becomes possible.
\item \textbf{Correct the documentation of ntuple 6} (F3), and if a real
transmitted fraction is needed, model the foil as a dielectric stack rather than
a metal.
\item \textbf{Reorder the polarisation assignment and restore the $2\pi$ azimuth
range} (F4).
\item \textbf{Add seed control} (section~\ref{sec:repro}) and \textbf{a
configurable output path} (F5), and \textbf{introduce a measured bulk
attenuation length} for the blue-WLS slab.
\item \textbf{Add at least one execution-level regression test.} The present
suite is source-level only and passed in full while the transport behaviour
described here was present. A short run asserting the sign of the
reflector-versus-no-reflector difference would have caught F1 immediately.
\end{enumerate}

% =============================================================================
\section*{Acknowledgments}
\addcontentsline{toc}{section}{Acknowledgments}

The code reading, build and run validation, diagnostic experiments and the
drafting of this note were carried out with the assistance of an AI coding
agent (Anthropic Claude, via Claude Code). The agent executed all builds, test
runs and simulation runs reported here and produced the first draft of this
document. All findings, in particular the mechanism proposed in
section~\ref{sec:f1}, are pending review and confirmation by the authors. This
note is an internal record and is not intended for publication.

\end{document}
